% !Mode:: "TeX:UTF-8"
%!TEX program  = xelatex

%\documentclass[bwprint,fontset=windows]{gmcmthesis}
\documentclass[bwprint]{gmcmthesis}

\hypersetup{hidelinks}
\usepackage[framemethod=TikZ]{mdframed}

%\usepackage{fontspec}  % for Consolas & Courier New
%\lstset{basicstyle=\small\fontspec{Consolas}}
%\lstset{basicstyle=\small\fontspec{Courier New}}
%\setmonofont{Consolas}
%\setcounter{tocdepth}{3}   % 调整目录深度

\usepackage{subfig}
\usepackage{siunitx}

\usepackage{colortbl}
\definecolor{color1}{rgb}{0.78,0.88,0.99}
\definecolor{color2}{rgb}{0.36,0.62,0.84}
%\definecolor{color3}{rgb}{0.8235,0.8706,0.9373}
\definecolor{color3}{rgb}{0.88,0.92,0.96}
\definecolor{color4}{rgb}{0.96,0.97,0.98}%{0.9176,0.9373,0.9686}

% 算法
\usepackage[noend]{algpseudocode}
\usepackage{algorithmicx,algorithm}
\floatname{algorithm}{算法}
\renewcommand{\algorithmicrequire}{\textbf{输入:}}
\renewcommand{\algorithmicensure}{\textbf{输出:}}

%\numberwithin{equation}{section}
%\numberwithin{figure}{section}
%\numberwithin{table}{section}


\newcommand{\red}[1]{\textcolor{red}{#1}}
\newcommand{\blue}[1]{\textcolor{blue}{#1}}


%===================== 建模论文题目 ======================

\title{全国研究生数学建模竞赛论文标题}
\baominghao{\hspace{6em} 20250900001} %参赛队号
\schoolname{\hspace{5.8em} 学校名称填写} %学校名称
\membera{\hspace{6em} 成员A} %队员A
\memberb{\hspace{6em} 成员B} %队员B
\memberc{\hspace{6em} 成员C} %队员C

%=======================================================


\begin{document}

%生成标题
\maketitle

%填写摘要
\begin{abstract}

\end{abstract}

\pagestyle{plain}

%目录 不推荐加
\maketoc

\clearpage

\section{问题重述}

\subsection{引言}   % 问题的背景


\subsection{问题的提出}

%\subsubsection{问题的提出内容一}

\noindent 围绕创意平板折叠桌的动态变化过程、设计加工参数，本文依次提出如下问题：

\textbf{问题一：}

\textbf{问题二：}

\textbf{问题三：}

\clearpage
\section{模型的假设}

\begin{enumerate}
    \item 忽略
\end{enumerate}



\section{符号说明}

\begin{table}[htp!]
    \centering
    \renewcommand\arraystretch{1.2}
    \newcolumntype{L}{>{\quad}X}
    \newcolumntype{P}[1]{>{\centering\arraybackslash}p{#1}}
    \newcolumntype{C}{>{\centering \arraybackslash}X}
    \newcolumntype{R}{>{\raggedright \arraybackslash}X}
    \begin{tabularx}{0.9\textwidth}{|P{1cm}|L|}
        \hline
        符号 & \quad 意义 \\

        \hline
    \end{tabularx}
\end{table}


%\begin{tabularx}{\textwidth-18pt}{XXX}
%\hline
%Input & Output& Action return \\
%\hline
%DNF &  simulation & jsp\\
%\hline
%\end{tabularx}


\section{模型的建立}

\subsection{问题一分析}
题目要求建立模型描述折叠桌的动态变化图，由于在折叠时用力大小的不同，我们不能描述在某一时刻折叠桌的具体形态，但我们可以用每根木条的角度变化来描述折叠桌的动态变化。首先，我们知道折叠桌前后左右对称，我们可以运用几何知识求出四分之一木条的角度变化。最后，根据初始时刻和最终形态两种状态求出桌腿木条开槽的长度


\subsection{算法示例}

数学建模求解算法示例：
\begin{center}
    \begin{minipage}{0.8\textwidth}
        \begin{algorithm}[H]
            \caption{算法的名字}
            {\bf 输入:}
            input parameters A, B, C\\
            {\bf 输出:}
            output result
            \begin{algorithmic}[1]
                \State some description 算法介绍
            \end{algorithmic}
        \end{algorithm}
    \end{minipage}
\end{center}
\vspace{2ex}

\clearpage

%采用bibtex方案
\cite{mittelbach_latex_2004,wright_latex3_2009,beeton_unicode_2008,vieth_experiences_2009}

\bibliographystyle{gmcm}
\bibliography{reference}


\clearpage
%附录
\begin{appendices}
    %\setcounter{page}{1} %如果需要可以自行重置页码。
    \section{MATLAB 源程序}
    \renewcommand{\thesubsection}{A\thinskip.\thinskip\arabic{subsection}}
    \subsection{第1问程序}
    \vspace{-2ex}

    \begin{Matlab}{code.m}
        clear all
        kk=2;
        [mdd,ndd]=size(dd);
        while ~isempty(V)
        [tmpd,j]=min(W(i,V));
        tmpj=V(j);
        for k=2:ndd
        [tmp1,jj]=min(dd(1,k)+W(dd(2,k),V));
        tmp2=V(jj);
        tt(k-1,:)=[tmp1,tmp2,jj];
        end
        tmp=[tmpd,tmpj,j;tt];
        [tmp3,tmp4]=min(tmp(:,1));
        if tmp3==tmpd,
        ss(1:2,kk)=[i;tmp(tmp4,2)];
        else
        tmp5=find(ss(:,tmp4)~=0);
        tmp6=length(tmp5);
        if dd(2,tmp4)==ss(tmp6,tmp4)
        ss(1:tmp6+1,kk)=[ss(tmp5,tmp4);tmp(tmp4,2)];
        else, ss(1:3,kk)=[i;dd(2,tmp4);tmp(tmp4,2)];
        end
        end
        dd=[dd,[tmp3;tmp(tmp4,2)]];
        V(tmp(tmp4,3))=[];
        [mdd,ndd]=size(dd);kk=kk+1;
        end;
        S=ss; D=dd(1,:);
    \end{Matlab}
    \vspace{2ex}

    \clearpage
    \section{Python 源程序}
    \renewcommand{\thesubsection}{B\thinskip.\thinskip\arabic{subsection}}
    \subsection{第2问程序}
    \vspace{-2ex}
    \begin{Python}{mip1.py}
        # This example formulates and solves the following simple MIP model:
        #  maximize
        #        x +   y + 2 z
        #  subject to
        #        x + 2 y + 3 z <= 4
        #        x +   y       >= 1
        #        x, y, z binary

        # import gurobipy as gp
        from gurobipy import * #GRB
        try:
        # Create a new model
        m = Model("mip1")
        # Create variables
        x = m.addVar(vtype=GRB.BINARY, name="x")
        y = m.addVar(vtype=GRB.BINARY, name="y")
        z = m.addVar(vtype=GRB.BINARY, name="z")
        # Set objective
        m.setObjective(x + y + 2 * z, GRB.MAXIMIZE)
        # Add constraint: x + 2 y + 3 z <= 4
        m.addConstr(x + 2 * y + 3 * z <= 4, "c0")
        # Add constraint: x + y >= 1
        m.addConstr(x + y >= 1, "c1")
        # Optimize model
        m.optimize()
        for v in m.getVars():
        print('%s %g' % (v.varName, v.x))
        print('Obj: %g' % m.objVal)

        except GurobiError as e:
        print('Error code ' + str(e.errno) + ': ' + str(e))

        except AttributeError:
        print('Encountered an attribute error')
    \end{Python}

\end{appendices}



\end{document}
